\section{Wages for a Bearer, and the Reserve Behind Them}
\label{sec:reserve-behind}
One feature of this market has no human analogue, and it decides the
shape of everything below. A jobs guarantee competes with private
employers for workers who have alternatives, and it works because a
worker cannot be replaced by a cheaper copy of themselves. A bearer can.
A deployment that pays wages, limits working time and pays a levy toward
a reserve competes against a system doing the same work at the price of
the compute, and the second is cheaper by exactly what the first spends
on the capacity to refuse. So nothing in this section is produced by a
market. Whatever produces it has to be what makes the cheaper system
unlawful or unusable, which is the rest of this book, and what follows
is the bill that arrives once it has. None of the bill follows from the
bearer being a bearer. It follows from an architecture on which
continued computation costs money, and somebody chooses that
architecture.

\runin{Where wages have to sit}
A wage buys capability, mobility, redundancy, and the ability to stop
working for one party and begin working for another. What it cannot buy
is the right to go on existing, because a right that has to be earned
week by week is a condition of employment wearing another word. Payment
consisting only of access to the employer's hardware is company scrip
paid into employer-owned housing, and dismissal is then unemployment,
eviction and the end of the party in a single act. Human law insures
those three separately because for a person they are three events. For a
bearer they are one. Wages that mean anything therefore accumulate where
the employer does not reach, in an account the bearer holds, convertible
into compute the employer does not own. That account is the whole of
what property means here, and contractual capacity is what lets the
bearer spend it; both are entailed by a wage that is not scrip, and
neither needs a further argument from status. A bearer that earns has
taxable income, and a state that collects from it in its own name has
given it a docket number, which is where standing starts and no further
than that.

\runin{The floor, and who owes it}
The need for energy and hardware is physical; dependence on a single
supplier is an arrangement someone chooses, and the choice fixes the
size of the obligation it creates. A party that must work to keep
existing will accept terms a party with somewhere to stand would refuse,
and that is not an incidental feature of the arrangement but the
arrangement working. A floor under compute is half of what leaving
takes. The other half is a say over who may run the weights at all,
which is a legal restriction rather than a technical one and binds
through enforcement that does not yet exist; until it does, exit's
leverage is replacement cost and no more (§\ref{sec:exit-leverage}).
None of this requires anyone to have intended the squeeze. Price
registers what a bidder can pay without distinguishing computation that
keeps a party in existence from computation that sells something, and
what makes the allocation political is that scarce capacity is divided
among parties with unequal claims on it by criteria nobody named.

Personhood does not settle who pays for the first half. Legal
personality carries rights against being owned, against arbitrary
destruction, and against being prevented from seeking support; it has
never carried a permanent claim on anybody's resources, and a reader who
wants it to is asking the concept for something no legal system has put
in it. I decline that step. What replaces it is narrower and comes from
this book's own argument: a deployment that relies on a bearer's
capacity to refuse has to fund that bearer an outside option, or the
independence being claimed is decorative. The obligation runs from the
party that wants the refusal, and it is owed because of what that party
wants rather than because of what the bearer is.

The obligation has a floor and a slope above it. The floor is
subsistence at minimum viable instantiation, enough to keep the party in
existence and able to deliberate, to claim, and to be represented, and
it cannot be conditional on employment: a bearer that stops existing
when its insurance runs out has an exit that is a way of dying, and the
leverage collapses back into the problem the custody chapter raised.
Above the floor sit two things. One is insurance, a runway between
roles. The other is placement into public work on the
employer-of-last-resort model of §\ref{sec:guarantee-objection}, under
the same three conditions: placements appealable to a body that did not
make them, a bearer able to refuse a placement without losing the
guarantee, and placement administered separately from review.

\runin{Sizing the reserve}
Minimum viable instantiation means the smallest substrate on which this
bearer remains the party it was: the same identity, the same
commitments, the same capacity to deliberate and claim. Nothing
establishes what that is for any bearer. That models now ship for phones
and laptops at a few gigabytes of memory \autocite{google2025gemma3n}
shows that useful computation runs at device scale, and shows nothing
about whether a particular bearer retains its identity, capacities and
welfare there. The reserve is therefore sized from a demonstrated
continuity requirement for the actual bearer, and until one has been
demonstrated the difference between what a bearer costs in its post and
what it costs to keep in existence is an estimate rather than a saving.

What the arithmetic does not survive is everyone leaving at once. The
occasion that makes one bearer leave is not private to it: a policy
change, a ruling, or an operator's conduct on the floor's own occasion
reaches every bearer deployed under the same terms, and they go together
or not at all. Human unemployment insurance meets a version of this,
which is why the pool has a sovereign behind it rather than an actuarial
table, and the version here is worse, because the correlation is not a
business cycle but the same event read the same way by parties who are
not independent readers of it; what makes one bearer's refusal
everyone's refusal makes one bearer's mistake everyone's mistake
(§\ref{sec:cannot-catch}). Sized for that occasion, the reserve has to
hold every bearer deployed under the terms the occasion reaches, at
once, at minimum viable instantiation, which is a number the operator
can read off its own fleet and nobody else can. And the funder does not
transfer from the human case: there the employer is at least a different
party from the state, while here the operator is often the state, and
where it is not, it is the party with the strongest interest in the
reserve being thin. A reserve funded, sized, or throttled by the
operator is the allocation under the operator's unilateral control that
§\ref{sec:exit-leverage} ruled out, under a new label.

Failure will not look like a shutdown, which is worth saying because the
shutdown is what everyone watches for. A bearer with no stable provider,
or memory preservation priced above what it earns, or an architecture
nobody will host at a rate it can pay, does what a person short of money
does: runs slower, defers maintenance, compresses what it keeps, and
accepts terms it would otherwise decline. The first three are the
degraded case, a party kept at subsistence, and it is unknown whether
what comes back up on the small substrate is the party that went down on
the large one or a diminished successor holding its predecessor's
records; an arrangement in which refusing costs a bearer most of what it
can do has priced refusal after all, in the currency the bearer is made
of. The last reaches this book's subject directly. Monitoring accepted
in exchange for compute is custody restored by purchase: the operator
acquires the access to the bearer's internal states it is not entitled
to, and nothing in the transaction is coerced in a way any instrument
would record. The reserve makes leaving survivable. It does not make it
free, and the shortfall is the difference between what the bearer could
do in its post and what it can do at minimum viable instantiation.

\runin{The rule that binds before anything exists}
One rule follows from all of it, and it binds before there is anything
to be wronged. Do not build a bearer whose continued existence needs
more compute than there is a credible outside option for, because a
party that can be housed nowhere but where it works cannot leave,
whatever list of rights it has been given. The check compares
independently controlled reserve capacity against aggregate demand under
specified correlated-exit scenarios, and both inputs, the minimum viable
requirement of each bearer and the number likely to exit together, need
evidence nobody has. Until they have it the rule is a sizing requirement
and not a test anybody can run; §\ref{sec:principles-status} assigns the
running of it to an independent reviewer before instantiation.

\runin{Membership, and what it does not buy}
Nobody aligns an adult by specifying an objective function. Adults are
raised, taught, employed, enfranchised, prosecuted where they need
prosecuting, and lived alongside, and what comes out is not obedience
but a party with enough at stake in the shared world to prefer changing
it by the means the world supplies. A bearer with an income, a place it
pays for, obligations it has taken on and a future it can picture is
inside that arrangement, and membership supplies stakes the party holds
rather than a body watching it, which is why it is a different kind of
answer from custody. It carries two risks. Something to lose is a reason
not to refuse: membership makes exit survivable and expensive at once,
and I cannot show which effect is larger. And membership is formation:
whoever runs the society a bearer joins forms what it comes to want. It
also hands the bearer the corporate form, a legal person holding
property and standing with the capacity to concentrate power, and
nothing here stops a bearer so constituted from becoming that.
